\chapter{Sprint 66 --- Cụm J Phần B + T: thương hiệu cấu hình được và tool Bản dịch (2026-10-04 $\to$ 10-07)}

\emph{Chapter này ghi lại hai phần còn lại của cụm J: Phần B cho đổi thương hiệu (tên, màu, logo, favicon, nền đăng nhập) ngay
trong Settings, và Phần T là tool Bản dịch trong backend (sửa, nhập/xuất, dịch máy). Phần V (Việt hoá source) ở chapter 79.
Phần O (Portal Vận hành) \textbf{tạm dừng theo ý BA} (07/10). Nhật ký từng phiên ở \texttt{docs/f-progress.md}; bài học ở
\texttt{docs/next-session-clusters-J.md} §3, §3b, §6.}

\section{Vì sao có cụm J}

Ngày 04/10 chủ dự án giao ba việc mới và chốt làm \textbf{trước} Issue List:
(1) thương hiệu cấu hình được trong Settings, để cùng source dùng cho thương hiệu khác;
(2) tool dịch đầy đủ trong backend, thay cho việc tự sửa và upload file \texttt{.po};
(3) portal Vận hành nhượng quyền trên backend \texttt{wujia\_franchise\_operations} sẵn có.
Ngày 05/10 cụm J gộp còn 7 phiên: J-T1 và J-T2 gộp, bỏ J-T3 (chuỗi trong code đổi thì xuất \texttt{.po} rồi restart), O3 và O4 gộp.
Phần V chen vào trước J-T4 sau khi máy quét đếm được 2\,659 chuỗi tiếng Việt viết cứng.

\section{Các phiên}

\begin{center}
\footnotesize
\begin{tabularx}{\textwidth}{>{\raggedright\arraybackslash}p{1.3cm}>{\raggedright\arraybackslash}p{2.6cm}>{\raggedright\arraybackslash}X>{\raggedright\arraybackslash}p{1.1cm}}
\toprule
Phiên & Commit & Nội dung & Ngày \\
\midrule
J-B1 & \texttt{de64704c} & \texttt{wujia\_core}: 5 field brand trên \texttt{res.company}, tab Settings ``Thương hiệu'', route ảnh, bộ màu & 04/10 \\
J-B2 & \texttt{12c2baa9} + \texttt{3aa974ac} & Áp brand vào portal (head, logo, CSS) và backend (favicon, tiêu đề tab) & 04/10 \\
J-T1+T2 & \texttt{ba08698e} & Module mới \texttt{wujia\_i18n}: danh mục chuỗi, màn sửa, Áp dụng không restart, giữ bản sửa qua \texttt{-u} & 05/10 \\
J-T4 & \texttt{ca198ace} & Nhập/xuất CSV kiểu file Thái, zip \texttt{.po}/\texttt{.pot}, CLI \texttt{scripts/i18n\_tool.py} & 07/10 \\
J-T5 & \texttt{78a68ea5} & Dịch máy DeepL: hàng đợi cron, trạng thái ``Dịch máy'', bảng thuật ngữ, Settings & 07/10 \\
J-T5b & --- & Chạy DeepL bằng key thật trên UAT (chỉ đọc + 1 lần kiểm hạn mức) & 07/10 \\
\bottomrule
\end{tabularx}
\end{center}

Trên UAT (kiểm chỉ đọc 07/10, khớp repo): \texttt{wujia\_core} 19.0.2.1.1 · \texttt{wujia\_portal\_layout} 19.0.60.4.0 ·
\texttt{wujia\_i18n} 19.0.1.3.0.

\section{Phần B --- thương hiệu}

\subsection*{Gì đã làm}
Quản trị vào Settings $\to$ ``Thương hiệu'' đặt tên, màu chính, logo PC/mobile, favicon và ảnh nền trang đăng nhập. Portal và
backend đổi theo ngay: tiêu đề tab ``\ldots{} · <tên brand>'', logo, favicon, màu nút/viền/nền nhạt. 25 câu ``Ngô Gia'' trong
QWeb và 15 chuỗi Python đọc tên từ cấu hình. Cấu hình mặc định giữ đúng giao diện cũ: đo 26 route $\times$ 3 khổ, 0 lệch pixel thật.

\begin{itemize}[nosep]
\item \texttt{wujia\_core}: \texttt{wj\_brand\_name}, \texttt{wj\_primary\_color}, \texttt{wj\_logo\_mobile}, \texttt{wj\_favicon},
  \texttt{wj\_login\_background}; logo PC dùng lại \texttt{res.company.logo}. \texttt{\_wj\_brand\_info()} có ormcache, ghi field brand
  thì xoá cache. Route \texttt{/wj/brand/<cid>/<kind>} public, có \texttt{?v=} nên cache \texttt{immutable}.
\item \texttt{tools/brand\_palette.py}: màu mặc định trả đúng bộ token BA với CSS rỗng; màu khác sinh dải HSL, nút CTA tự làm đậm tới
  tương phản chữ trắng $\geq$ 4.5:1; \texttt{retint()} xoay sắc 12 token xanh còn sót (gradient, nền ô).
\item \texttt{wujia\_portal\_layout}: biến QWeb \texttt{wj\_brand} chèn ở \texttt{ir.qweb.\_prepare\_environment} (route portal không
  \texttt{website=True}); logo qua URL thay base64 nên HTML 26 trang nhẹ đi 36\% (\texttt{/portal/login} 46 KB $\to$ 13 KB).
\item Câu có tên brand: mẫu \texttt{<t t-set="wj\_txt">\ldots\{brand\}\ldots</t>} + \texttt{.replace()} để vẫn là một term dịch được.
\end{itemize}

\subsection*{Lý do / trade-off}
Không dùng \texttt{res.company.primary\_color} sẵn có vì đó là màu báo cáo PDF, wizard Document Layout tự đổi theo logo. Tên brand
mặc định ``Ngô Gia'' tự điền khi cài hoặc nâng cấp (data \texttt{noupdate} cho \texttt{-i} cộng migration cho \texttt{-u}), không
ghi đè tên đã đặt. Biểu đồ màn Khảo sát (code anh Thái) vẫn màu cứng, để nhóm đó chỉnh.

\section{Phần T --- tool Bản dịch}

\subsection*{Gì đã làm}
Người dịch có app ``Translation Tool'' trong backend. Bấm Quét để gom mọi chuỗi của module team theo ngôn ngữ, sửa ngay trên màn,
bấm Áp dụng thì nhãn, menu, view, QWeb portal đổi \textbf{ngay, không restart}. Bản sửa tay giữ được qua các lần nâng cấp module.
Có thể nhập/xuất CSV (đúng định dạng file anh Thái đang dùng) và xuất zip \texttt{.po}/\texttt{.pot} cho chuỗi trong code. Dịch máy
DeepL: chọn ngôn ngữ và module, hệ thống dịch hàng loạt theo lô, áp ngay, người rà sau bấm ``Mark reviewed''.

\begin{itemize}[nosep]
\item Model: \texttt{wujia.i18n.term} (chuỗi) $\times$ \texttt{wujia.i18n.value} (bản dịch theo ngôn ngữ, trạng thái
  \texttt{synced | override | machine | missing}), \texttt{wujia.i18n.coverage} (view SQL độ phủ), \texttt{wujia.i18n.glossary}.
\item Áp dụng = \texttt{TranslationImporter} ghi đè + xoá đúng loại cache. Override \texttt{ir.module.module.\_update\_translations}
  áp lại bản sửa tay sau mỗi \texttt{-u}.
\item \texttt{tools/po\_writer.py} (không import odoo) dùng chung cho wizard, CLI \texttt{scripts/i18n\_tool.py} và
  \texttt{scripts/sync\_translations.py}. \texttt{tools/mt\_deepl.py}: bảo vệ placeholder/thẻ HTML, mất placeholder thì loại câu.
\item Cron hàng đợi dịch máy: lô 50 câu/30\,000 ký tự, đọc lại trạng thái trước khi ghi (người sửa giữa chừng thắng), hết hạn mức
  hoặc sai key thì dừng và giữ hàng đợi.
\end{itemize}

\subsection*{Lý do / trade-off}
\begin{itemize}[nosep]
\item Bỏ J-T3 (lớp phủ chuỗi Python/JS không restart): chuỗi trong code chỉ có hiệu lực sau khi xuất \texttt{.po}, commit và restart.
  Đơn giản hơn, không rủi ro trên nhiều worker.
\item Nhập CSV và dịch máy \textbf{không đè bản sửa tay}; dòng chỉ khớp theo câu nguồn chỉ điền chỗ trống. Đo thử chế độ ``đè hết''
  làm đổi 90 + 162 bản vi\_VN đã chốt ở Phần V.
\item Dịch máy ``áp ngay, rà sau'' (chủ dự án chốt): nhanh có giao diện th/zh, đổi lại có thể có câu gượng tới khi được rà.
\end{itemize}

\subsection*{Số đo chính}
Quét 29 module $\times$ 3 ngôn ngữ: 4\,839 chuỗi trong 3,1 s nên chạy đồng bộ, không cần cron. Suite \texttt{wujia\_i18n} 55/0
(DB trắng và DB copy), mutation T1 7/7, T4 6/6, T5 8/8. Nhập glossary 2\,195 dòng trong 0,2 s. Smoke DeepL thật trên UAT:
th\_TH 1\,774 và zh\_CN 1\,850 chuỗi dịch máy, 2 câu bị bộ kiểm chặn đúng, portal th 6 route $\times$ 2 khổ: 0 tràn, 0 lỗi JS,
chữ Thái hiện ngay không restart.

\section{Tổng kết cụm J}

\begin{itemize}[nosep]
\item \textbf{Phần B} xong: đổi thương hiệu không cần sửa code.
\item \textbf{Phần V} xong (chapter 79): câu gốc source team là tiếng Anh, vi\_VN qua \texttt{.po}.
\item \textbf{Phần T} xong: dịch tay, nhập/xuất, dịch máy đều chạy trên UAT.
\item \textbf{Phần O} (J-O0 \ldots{} ★JR) \textbf{tạm dừng theo ý BA}. Câu hỏi treo: chấm công/nghỉ phép, ai duyệt ca, Staff xem gì,
  gắn nhân viên với tài khoản portal, spec Model Field mục G lệch backend, Figma.
\end{itemize}
Việc kế: Issue List cụm I mở rộng (I1 \ldots{} I10 + ★IR, 15 issue), rồi cụm H. Lộ trình ở \texttt{docs/next-session-clusters-H.md}.

\section{Bài học}

\begin{itemize}[nosep]
\item zsh không tách biến không ngoặc: một lệnh copy DB chạy \texttt{rm -rf data/filestore/} và mất filestore local. Lệnh nhiều tham số
  viết bằng file \texttt{\#!/bin/bash} + \texttt{set -u}; không \texttt{rm -rf} đường dẫn ghép biến khi chưa kiểm biến khác rỗng.
  Khôi phục: \texttt{scripts/dev/filestore\_pack.py} + \texttt{filestore\_restore.py}.
\item \texttt{res.lang.\_activate\_lang} chỉ bật ngôn ngữ, không nạp \texttt{.po}; phải gọi \texttt{\_update\_translations}.
\item \texttt{<function>} trong data \texttt{noupdate} chỉ chạy lúc \texttt{-i}; cần migration cho \texttt{-u}.
\item Glossary dùng chung khớp theo câu nguồn sẽ đè bản dịch riêng của module, nên mọi đường nhập/dịch máy chỉ điền chỗ trống.
\item Người dịch không phải admin không đọc được \texttt{ir.module.module}: cần quyền chỉ đọc, nếu không mọi wizard có ô Module lỗi.
\item \texttt{\_commit\_progress} ngoài cron là commit thật, chỉ gọi khi context có \texttt{cron\_id}.
\item Test chống lùi phải chọn dữ liệu bắt được đột biến: chuỗi code (DB không lưu bản dịch) chứ không phải nhãn; câu nguồn duy nhất
  chứ không phải câu lặp ở menu/action/view.
\end{itemize}

\section{Nợ để lại}

\begin{itemize}[nosep]
\item Xuất \texttt{.po} th/zh cho chuỗi trong code (th còn 276 chuỗi \texttt{\_()} chờ), commit rồi \texttt{-u wujia\_i18n} + restart;
  sửa tay 2 câu dịch máy bị chặn (zh exam ``HTML tags changed'', th sale \texttt{On \%s}).
\item Gửi BA 1\,361 cặp câu và gửi anh Thái 284 chuỗi + lỗi \texttt{.po} \texttt{web\_survey\_ui} (\texttt{docs/i18n-review/}).
\item Biểu đồ Khảo sát chưa theo màu brand (code anh Thái).
\item Phần O chờ BA.
\end{itemize}
